\section{Giới thiệu}
\label{sec:introduction}

\subsection{Đặt vấn đề}
Trong cuộc sống hiện đại, đặc biệt là đối với sinh viên, nhân viên văn phòng và các cặp đôi, câu hỏi "Hôm nay ăn gì?" đã phát triển từ một câu hỏi đơn giản thành một nguồn gây căng thẳng hàng ngày. Vấn đề này biểu hiện qua ba điểm yếu (pain points) chính:
\begin{itemize}
    \item \textbf{Sự mệt mỏi khi ra quyết định (Decision Fatigue):} Trung bình, mỗi cá nhân dành từ 15 đến 20 phút để lướt các ứng dụng giao đồ ăn một cách vô định mà không đưa ra được lựa chọn cuối cùng, thường dẫn đến việc chọn một bữa ăn ngẫu nhiên, không hài lòng hoặc bỏ bữa hoàn toàn.
    \item \textbf{Quá tải lựa chọn (Choice Overload):} Các nền tảng hiện nay hiển thị hàng trăm tùy chọn nhà hàng cùng lúc dưới dạng danh sách. Lượng thông tin khổng lồ này gây ra tình trạng quá tải nhận thức, khiến người dùng khó đưa ra quyết định và ít hài lòng hơn với lựa chọn cuối cùng của mình.
    \item \textbf{Mâu thuẫn nhóm (Social Friction):} Việc quyết định địa điểm ăn uống là một nguồn gây mâu thuẫn phổ biến cho các nhóm, thường dẫn đến sự bực bội hoặc những thỏa hiệp không thoải mái.
\end{itemize}

\subsection{Mục tiêu và Mục đích}
Mục tiêu chính của dự án này là phát triển \textbf{Tinner} (một từ ghép của "Tinder" và "Dinner"), một ứng dụng di động được thiết kế để giải quyết vấn đề "Tê liệt phân tích" (Analysis Paralysis) trong việc lựa chọn thực phẩm. Với khẩu hiệu \textit{"Swipe To Bite"}, Tinner đơn giản hóa quy trình ra quyết định.

Các mục đích cốt lõi của ứng dụng bao gồm:
\begin{itemize}
    \item \textbf{Hợp lý hóa các quyết định:} Giúp người dùng chuyển từ việc mở ứng dụng sang quyết định một bữa ăn trong vòng chưa đầy 5 phút.
    \item \textbf{Trải nghiệm ưu tiên hình ảnh (Visual-First Experience):} Cung cấp một "Giao diện chế độ xem đơn" (Single-View Interface) hiển thị hình ảnh món ăn chất lượng cao tại một thời điểm, loại bỏ các yếu tố gây xao nhãng và kích thích sự thèm ăn.
    \item \textbf{Khám phá thông qua Gamification:} Triển khai cơ chế quẹt thẻ (Vuốt sang trái để bỏ qua, Vuốt sang phải để lưu/thích) giúp biến một công việc tẻ nhạt thành một hoạt động hấp dẫn.
    \item \textbf{Truy cập thông tin liền mạch:} Khi thích một món ăn, ứng dụng ngay lập tức hiển thị Thẻ thông tin nhà hàng với các chi tiết quan trọng (địa chỉ, giờ mở cửa, tích hợp bản đồ và đánh giá) để hoàn tất quyết định.
    \item \textbf{Quản lý hành động:} Cho phép người dùng duy trì danh sách "Muốn thử" cá nhân để xem xét, so sánh và điều hướng đến các nhà hàng đã chọn một cách hiệu quả.
\end{itemize}

\subsection{Nhu cầu thực tế và Tác động tiềm năng}
Nghiên cứu thị trường toàn diện chỉ ra một khoảng trống đáng kể trong hệ sinh thái ứng dụng liên quan đến thực phẩm hiện nay. Trong khi các siêu ứng dụng giao hàng vượt trội về hậu cần, các nền tảng mạng xã hội vượt trội về cảm hứng hình ảnh và dịch vụ bản đồ vượt trội về dữ liệu vị trí, thì không có mô hình lai (hybrid model) ưu việt nào được tối ưu hóa thuần túy cho việc khám phá nhà hàng nhanh chóng, thú vị và có cấu trúc.

Tinner giải quyết nhu cầu thực tế này bằng cách định vị mình là một "Công cụ ra quyết định" (Decision Engine) hơn là một trung gian giao dịch. Bằng cách tập trung vào giai đoạn khám phá và kết hợp các tính năng như Bộ lọc thông minh (Smart Filters) và Bộ sưu tập cá nhân (Personal Collections), ứng dụng giúp giảm tải nhận thức và thu hẹp khoảng cách giữa thông tin có cấu trúc và trải nghiệm người dùng hấp dẫn.

Tác động tiềm năng của Tinner là rất lớn, đặc biệt là đối với Gen Z và Millennials. Nó hứa hẹn sẽ giảm đáng kể thời gian lãng phí cho các quyết định ăn uống hàng ngày, giảm bớt mâu thuẫn xã hội liên quan đến lựa chọn của nhóm và cung cấp một giải pháp trực quan và thú vị cho một vấn đề nan giải hàng ngày.
